% Original Prism writing draft for European Journal of Medicinal Chemistry.
% Distribution-provided elsarticle and elsarticle-num remain unchanged.
% All demonstration values and graphics are synthetic, not research findings.
\documentclass[a4paper,preprint,12pt]{elsarticle}
\usepackage{amsmath,amssymb,graphicx,array,booktabs,tabularx,siunitx,newfloat}
\usepackage[version=4]{mhchem}
\usepackage[hidelinks]{hyperref}
\providecommand{\doi}[1]{\href{https://doi.org/#1}{\nolinkurl{#1}}}
\DeclareFloatingEnvironment[fileext=los,name=Scheme,placement=htbp]{scheme}
\biboptions{numbers,sort&compress}
\setcitestyle{square}
\sisetup{uncertainty-mode=separate}
\setlength{\emergencystretch}{2em}
\raggedbottom
\clubpenalty=10000
\widowpenalty=10000
\displaywidowpenalty=10000
\renewcommand{\topfraction}{0.9}
\renewcommand{\textfraction}{0.08}
\renewcommand{\floatpagefraction}{0.8}
% Keep draft float pages compact instead of stretching space between graphics.
\makeatletter
\setlength{\@fptop}{0pt}
\setlength{\@fpsep}{18pt plus 2pt minus 2pt}
\setlength{\@fpbot}{0pt plus 1fil}
\makeatother
\journal{European Journal of Medicinal Chemistry}

\begin{document}
\begin{frontmatter}
\title{Manuscript title}
\author[inst1]{First Author}
\author[inst1,inst2]{Second Author\texorpdfstring{\corref{corresponding}}{}}
\cortext[corresponding]{Corresponding author.}
\ead{second.author@example.org}
\address[inst1]{Department, First Institution, City, Country}
\address[inst2]{Department, Second Institution, City, Country}

\begin{abstract}
Replace this paragraph with a concise, self-contained account of the
question, approach, verified findings and conclusions. Follow the current
journal instructions for the selected article type. These resources are
synthetic layout demonstrations: no compounds were prepared or tested and
no scientific result is reported. Replace every demonstration before use.
\end{abstract}
\begin{keyword}
keyword one \sep keyword two \sep keyword three
\end{keyword}
\end{frontmatter}

\section{Introduction}\label{sec:introduction}
Explain the medicinal chemistry question, rationale and contribution.
Discuss the relevant target or mechanism and the evidence supporting the
study design. The public citations here exercise reference ordering only,
not medicinal chemistry evidence~\cite{wilkinson2016fair,peng2011reproducible}.
Retain them only if relevant to the actual manuscript. Replace every
instruction, invented number and framed graphic before submission.

\section{Results and Discussion}\label{sec:results}
\subsection{Design and chemistry}\label{sec:chemistry}
Report the authors' verified design and chemical work; distinguish measured
findings from hypotheses or computational predictions.

Scheme~\ref{sch:sequence} demonstrates an editable vector reaction-layout
placeholder. Use the authors' original structures and preserve compound
identities consistently in the text, schemes and tables.
\begin{scheme}[htbp]
\centering
\fbox{\parbox[c][24mm][c]{0.86\linewidth}{\centering
\ce{A ->[step 1] B ->[step 2] C}\\[7pt]
Replace A, B and C with original structures and verified conditions.}}
\caption{\textbf{Synthetic scheme layout.} Letters stand for layout slots,
not chemical substances. No reaction, yield, synthesis route or experiment
is reported. Replace the drawing, conditions and legend with the authors'
verified work, or remove this demonstration.}
\label{sch:sequence}
\end{scheme}

Figure~\ref{fig:overview} demonstrates a two-panel layout. Long legends
belong to the page layout: preserve readable panel labels and account for
the full caption height when choosing the graphic size.
\begin{figure}[htbp]
\centering
\begin{minipage}[t]{0.47\linewidth}
\centering\textbf{a}\par\smallskip
\fbox{\parbox[c][27mm][c]{0.90\linewidth}{\centering
Original structure slot\\[5pt]Replace with a vector graphic}}
\end{minipage}\hfill
\begin{minipage}[t]{0.47\linewidth}
\centering\textbf{b}\par\smallskip
\fbox{\parbox[c][27mm][c]{0.90\linewidth}{\centering
Evidence graphic slot\\[5pt]Keep final-size text readable}}
\end{minipage}
\caption{\textbf{Synthetic two-panel example with an extended legend.}
(a) An original structure or concept diagram may occupy this slot.
(b) An independently prepared evidence display may occupy the second slot.
These framed regions contain no data. Define the actual symbols, units,
uncertainty and relevant conditions in the replacement legend. Preserve
panel identities across text references and graphic labels; keep axes,
legends and chemical labels readable at the final printed size. Remove
instructions and unused panels before submission.}
\label{fig:overview}
\end{figure}

Table~\ref{tab:sar} illustrates a structure--activity relationship (SAR)
table with bold compound identifiers, chemical formulae and aligned errors.
The entries are invented formatting values, not assay observations.
\begin{table}[htbp]
\centering\small
\caption{Synthetic SAR-style values for column-alignment practice only.}
\label{tab:sar}
\begin{tabularx}{\linewidth}{@{}c >{\raggedright\arraybackslash}X
 S[table-format=2.1(2)] S[table-format=2.1(2)]@{}}
\toprule
Compound & Substituent slot & {\shortstack{Illustrative endpoint 1\\/ nM}}
 & {\shortstack{Illustrative endpoint 2\\/ nM}}\\
\midrule
\textbf{1} & \ce{CH3} & 12.4 +- 0.6 & 31.0 +- 1.4\\
\textbf{2} & \ce{C2H5} & 8.1 +- 0.4 & 27.6 +- 1.2\\
\textbf{3} & \ce{F} & 15.2 +- 0.8 & 19.4 +- 0.9\\
\textbf{4} & Replace with an actual group & \multicolumn{1}{c}{---}
 & \multicolumn{1}{c}{---}\\
\bottomrule
\end{tabularx}
\par\smallskip\footnotesize\raggedright
No assay was performed. Replace all values and define the actual endpoint,
units, error measure and replicate counts. A dash is a missing-value slot,
not a zero. Formulae such as \ce{CO2} demonstrate chemical typography only.
\end{table}

Equations~\eqref{eq:ratio} and~\eqref{eq:difference} demonstrate alignment
and cross-references. Define actual quantities before using any equation.
\begin{align}
 q_i &= \frac{x_i}{x_{\mathrm{ref}}},\label{eq:ratio}\\
 \Delta q &= q_{\mathrm{A}}-q_{\mathrm{B}}.\label{eq:difference}
\end{align}
These expressions are layout examples, not a validated medicinal chemistry
metric or evidence of a scientific result.

\subsection{Biological evaluation and interpretation}\label{sec:biology}
Describe actual assays, controls, statistics and limitations. Relate each
conclusion to the reported evidence, distinguishing measurements from
mechanistic interpretation. Do not infer experimental validation from a
computed result alone.

\section{Conclusion}\label{sec:conclusion}
Summarize conclusions supported by the study and state the limits of their
interpretation. Identify next steps without presenting them as completed.

\section{Experimental}\label{sec:experimental}
\subsection{General methods and characterization}
Supply the actual materials, instruments, analytical methods and software
versions. Report measured purity and structure characterization with their
supporting evidence; do not retain an unverified purity or identity claim.
Describe relevant ethics approvals and unusual hazards when applicable.

\subsection{Synthesis and biological methods}
Replace this section with the procedures actually performed, their verified
conditions, characterization and reproducible assay methods. Determine the
appropriate division between main text and supporting files using the
current journal guidelines. This starter contains no experimental result
or claim that a synthesis, assay or validation was performed.

\section*{Data availability}
Replace with the authors' accurate availability statement, repository
identifiers where applicable, and any access restrictions.
\section*{Author contributions}
Describe actual contributions using the journal's current requirements.
\section*{Declaration of competing interest}
Replace with the authors' confirmed declaration.
\section*{Acknowledgements and funding}
List confirmed acknowledgements and funding, or remove if not applicable.
\section*{Supplementary material}
Describe the actual supporting files. The included \texttt{supplementary.tex}
compiles independently with S-numbered pages and displays. A separate
\texttt{graphical\_abstract.tex} and editable \texttt{highlights.txt} are
optional drafting aids; confirm their submission requirements with the
journal before using them.

% The two public references test mechanics only; replace as appropriate.
\bibliographystyle{elsarticle-num}
\bibliography{references}
\end{document}
